\documentclass[11pt]{article}
\usepackage[T1]{fontenc}
\usepackage{lmodern}
\usepackage{amsmath, amssymb, amsthm}
\usepackage[margin=1.2in]{geometry}

\newtheorem{theorem}{Theorem}
\newtheorem{lemma}[theorem]{Lemma}
\newtheorem{corollary}[theorem]{Corollary}
\newtheorem{definition}[theorem]{Definition}
\newtheorem{proposition}[theorem]{Proposition}

\title{\textbf{Frozen-Jacobian Divergence Analysis}}
\author{}
\date{\vspace{-5ex}}

\begin{document}
\maketitle

\section{Setup}

Let $\phi$ be a self-concordant barrier on a cone $K$ with barrier
parameter $\nu$.  Let $z \in \mathrm{int}(K)$ be the current iterate
and $z_0 \in \mathrm{int}(K)$ the reference point at which the Hessian
was last factored.  The constraint map is $z = Ax + b$ with
$A \in \mathbb{R}^{m \times n}$.

Let $z_*$ be a fixed target point and define the Jeffreys divergence
\begin{equation}
    V(z) = \langle z - z_*, \nabla\phi(z) - \nabla\phi(z_*) \rangle.
\end{equation}
Define the two kinetic energies
\begin{equation}
    E = \|\dot{z}\|_z^2 = \dot{z}^T H(z)\dot{z}, \qquad
    E_0 = \|\dot{z}\|_{z_0}^2 = \dot{z}^T H(z_0)\dot{z},
\end{equation}
the Hessian distortion ratio $\sigma = E_0 / E$, and the infinity-norm
energy $E_\infty = \|\dot{z}\|_{z,\infty}^2$ (i.e., the squared
maximum eigenvalue of $\dot{z}$ in the local metric at $z$; for
nonneg, $E_\infty = d_\infty^2 = \max_i(\dot{z}_i/z_i)^2$).

\section{Frozen-Jacobian KKT Orthogonality}

The frozen Newton direction $\dot{z} \in \mathrm{range}(A)$ is
determined by the KKT system with Gram $A^T H(z_0) A$.  The resulting
direction satisfies:
\begin{equation} \label{eq:frozen_kkt}
    \nabla\phi(z) + H(z_0)\dot{z} - \nabla\phi(z_*) \perp \mathrm{range}(A).
\end{equation}
Since both $\dot{z} \in \mathrm{range}(A)$ and
$z - z_* \in \mathrm{range}(A)$, substituting each into
\eqref{eq:frozen_kkt} gives:
\begin{align}
    \langle \dot{z},\; \nabla\phi(z) - \nabla\phi(z_*) \rangle
        &= -E_0, \label{eq:orth_zdot} \\
    V &= -\langle z - z_*,\; H(z_0)\dot{z} \rangle. \label{eq:orth_v}
\end{align}

\section{First Derivative of the Divergence}

\begin{lemma}[Frozen $\dot{V}$ Identity]
\label{lem:vdot_frozen}
Under the frozen-Jacobian Newton direction,
\begin{equation} \label{eq:vdot_frozen}
    \dot{V} = -(E_0 + V) + \epsilon,
\end{equation}
where $\epsilon = \langle z - z_*,\; \Delta H\, \dot{z} \rangle$
with $\Delta H = H(z) - H(z_0)$.
\end{lemma}

\begin{proof}
By the chain rule:
$\dot{V} = \langle \dot{z}, \nabla\phi(z) - \nabla\phi(z_*) \rangle
         + \langle z - z_*, H(z)\dot{z} \rangle$.
Substituting \eqref{eq:orth_zdot} and decomposing $H(z) = H(z_0) + \Delta H$:
$\dot{V} = -E_0 + \langle z-z_*, H(z_0)\dot{z}\rangle + \epsilon
         = -E_0 - V + \epsilon$.
\end{proof}

\begin{lemma}[Error Bound]
\label{lem:epsilon_bound}
If $\rho = \|z - z_0\|_z < 1$, then
\begin{equation}
    |\epsilon| \le \frac{\rho}{1-\rho}\,(V + \sqrt{V})\,\sqrt{E}.
\end{equation}
\end{lemma}

\begin{proof}
Cauchy--Schwarz in the $H(z)$-metric:
$|\epsilon| \le \|z-z_*\|_z \cdot \|H(z)^{-1}\Delta H\,\dot{z}\|_z$.
The first factor satisfies $\|z-z_*\|_z \le V + \sqrt{V}$
(self-concordant bound).  The second satisfies
$\|H(z)^{-1}\Delta H\,\dot{z}\|_z \le \frac{\rho}{1-\rho}\sqrt{E}$
(integral metric comparison).
\end{proof}

\section{Second Derivative: General Case}

\begin{lemma}[$\ddot{V}$ Exact Identity]
\label{lem:vddot}
The geodesic follows the true metric $H(z)$ regardless of which
Hessian was used in the Newton solve:
\begin{equation} \label{eq:vddot_exact}
    \ddot{V} = 2E - \langle \ddot{z}, U \rangle_z,
\end{equation}
where $U = (z - z_*) - H(z)^{-1}(\nabla\phi(z) - \nabla\phi(z_*))$.
\end{lemma}

\begin{lemma}[General Curvature Bounds]
\label{lem:curvature_general}
Under the kinematic stability axiom $\|\ddot{z}\|_z \le E$:
\begin{enumerate}
    \item \textbf{Global:}
    $\ddot{V} \le 2E(1 + \sqrt{V} + V)$.
    \item \textbf{Near-field} ($E < 1/4$):
    $\ddot{V} \le 2E(1 + V)$.
\end{enumerate}
\end{lemma}

\begin{proof}
From \eqref{eq:vddot_exact} and Cauchy--Schwarz:
$\ddot{V} \le 2E + E\|U\|_z$.  The global bound uses
$\|U\|_z \le 2V + 2\sqrt{V}$ (triangle inequality).  The near-field
bound uses $\|U\|_z \le 2V$ (coupled integral bound when
$\|z-z_*\|_z < 1$).
\end{proof}

\section{Second Derivative: Symmetric Cones}

For symmetric cones (nonneg, SOC, PSD), the barrier $\phi$ defines a
Euclidean Jordan algebra.  In the isotropic frame at $z$, the Newton
direction has spectral components $d_i = \tilde{\dot{z}}_i$ with
$E = \sum_i d_i^2$, $E_\infty = \max_i d_i^2 = d_\infty^2$.  The
relative stretch vector is $w_i = (\text{eigenvalues of } P(z)^{-1/2}z_*)_i$.

\begin{lemma}[Symmetric Cone Exact Formula]
\label{lem:sym_exact}
For symmetric cones, the inner product evaluates exactly to:
\begin{equation} \label{eq:sym_exact}
    \langle \ddot{z}, U \rangle_z = \sum_i d_i^2\,(2 - w_i - w_i^{-1}).
\end{equation}
Since $w_i + w_i^{-1} \ge 2$ by AM-GM, each term is non-positive, so
$\langle \ddot{z}, U \rangle_z \le 0$ and:
\begin{equation}
    \ddot{V} = 2E + \sum_i d_i^2\,(w_i + w_i^{-1} - 2) \ge 2E.
\end{equation}
\end{lemma}

\begin{proof}
In the isotropic frame, the geodesic acceleration satisfies
$\tilde{\ddot{z}}_i = d_i^2/w_i$ (from the Jordan algebra geodesic
equation).  The residual has spectral components
$\tilde{U}_i = z_i(2 - w_i - w_i^{-1})$.  The local inner product
$\langle \ddot{z}, U \rangle_z = \sum_i \tilde{\ddot{z}}_i
\tilde{U}_i H_{ii}$ reduces to \eqref{eq:sym_exact} after
cancellation of the metric factors.
\end{proof}

\begin{theorem}[Symmetric Cone Curvature Bound]
\label{thm:sym_curvature}
For symmetric cones:
\begin{equation} \label{eq:sym_curvature}
    \ddot{V} \le 2E + d_\infty^2\, V.
\end{equation}
\end{theorem}

\begin{proof}
From \eqref{eq:sym_exact}, define $f_i = w_i + w_i^{-1} - 2 \ge 0$.
Then $\ddot{V} = 2E + \sum_i d_i^2 f_i$.  Apply H\"older's inequality
($\ell^\infty \cdot \ell^1$):
\begin{equation}
    \sum_i d_i^2 f_i \le \max_i(d_i^2) \cdot \sum_i f_i
    = d_\infty^2 \cdot V,
\end{equation}
where the last equality uses $V = \sum_i (w_i + w_i^{-1} - 2) = \sum_i f_i$
(the spectral decomposition of the Jeffreys divergence for symmetric
cones).
\end{proof}

\begin{corollary}
The symmetric cone bound \eqref{eq:sym_curvature} is strictly tighter
than the general bound (Lemma~\ref{lem:curvature_general}) whenever
$d_\infty^2 < 2E$, i.e., the direction is not concentrated on a single
eigenvalue.  Since $E = \sum_i d_i^2 \ge d_\infty^2$ with equality
iff all mass is on one component, the symmetric cone bound replaces
$2EV$ (from the general near-field) with $d_\infty^2 V$---a factor
of $E/d_\infty^2 = \nu_{\text{eff}}$ improvement, where
$\nu_{\text{eff}}$ is the effective number of active eigenvalues.
\end{corollary}

\section{Frozen-J Descent on Symmetric Cones}

We now combine the frozen $\dot{V}$ identity
(Lemma~\ref{lem:vdot_frozen}) with the symmetric cone curvature bound
(Theorem~\ref{thm:sym_curvature}).  The Taylor surrogate becomes:
\begin{equation} \label{eq:taylor_sym}
    V(h) \le V - h(E_0 + V - |\epsilon|) + \tfrac{1}{2}h^2(2E + d_\infty^2 V).
\end{equation}

\begin{theorem}[Frozen-J Descent, Symmetric Cone, Far-Field]
\label{thm:sym_farfield}
Suppose $E \ge 1$ and $|\epsilon| \le E_0/2$.  The step size
\begin{equation}
    h_f = \frac{E_0}{2(2E + d_\infty^2 V)}
\end{equation}
guarantees strict descent with
\begin{equation}
    \Delta V \le -\frac{E_0^2}{8(2E + d_\infty^2 V)}.
\end{equation}
\end{theorem}

\begin{proof}
From \eqref{eq:taylor_sym}, descent holds when
$h < D/C$ where $D = E_0 + V - |\epsilon| \ge E_0/2$ and
$C = \frac{1}{2}(2E + d_\infty^2 V)$.  Setting $h_f = D/(2C) \ge
E_0/(4C) = E_0/(2(2E + d_\infty^2 V))$:
\begin{equation}
    \Delta V \le -h_f D + h_f^2 C = -\frac{D^2}{4C}
    \le -\frac{E_0^2/4}{2(2E + d_\infty^2 V)/2}
    = -\frac{E_0^2}{8(2E + d_\infty^2 V)}.
\end{equation}
Geodesic feasibility ($h_f \le \alpha$): since $\alpha = 2/d_\infty^2$
in the far-field, $h_f \le \alpha$ requires
$E_0 d_\infty^2 \le 4(2E + d_\infty^2 V)$, which holds because
$E_0 \le 4E$ and $d_\infty^2 V \ge 0$.
\end{proof}

\begin{theorem}[Frozen-J Descent, Symmetric Cone, Near-Field]
\label{thm:sym_nearfield}
Suppose $d_\infty^2 < 2$ (so $\alpha = 1$) and $|\epsilon| \le E_0/2$.
The step size $h_f = 1$ guarantees strict descent with
\begin{equation}
    V(1) \le \tfrac{1}{2}d_\infty^2 V + (E - E_0) + |\epsilon|.
\end{equation}
When $z_0 = z$ (unfrozen): $V(1) \le \frac{1}{2}d_\infty^2 V$
(pure quadratic convergence with rate $d_\infty^2/2$ instead of $E$).
\end{theorem}

\begin{proof}
Substituting $h = 1$ into \eqref{eq:taylor_sym}:
\begin{align}
    V(1) &\le V - (E_0 + V - |\epsilon|) + \tfrac{1}{2}(2E + d_\infty^2 V) \nonumber \\
         &= V - E_0 - V + |\epsilon| + E + \tfrac{1}{2}d_\infty^2 V \nonumber \\
         &= (E - E_0) + \tfrac{1}{2}d_\infty^2 V + |\epsilon|.
\end{align}
For the unfrozen case ($z_0 = z$): $E_0 = E$, $\epsilon = 0$, giving
$V(1) \le \frac{1}{2}d_\infty^2 V$.

For descent, we need $\frac{1}{2}d_\infty^2 V + |E-E_0| + |\epsilon| < V$,
i.e., $(1 - \frac{1}{2}d_\infty^2)V > |E-E_0| + |\epsilon|$.  Since
$d_\infty^2 < 2$, the left side is positive, and the condition holds
when the frozen perturbations are small relative to $V$.
\end{proof}

\begin{corollary}[Improvement over General Near-Field]
\label{cor:improvement}
In the unfrozen near-field, the general bound gives $V(1) \le EV$
(Theorem in \texttt{centering.tex}), while the symmetric cone bound
gives $V(1) \le \frac{1}{2}d_\infty^2 V$.  The improvement factor is
$2E/d_\infty^2$.  For well-distributed directions ($E = \nu d_\infty^2/m$
where $m$ is the number of eigenvalues), this is $2\nu/m$---significant
for large cones.
\end{corollary}

\section{Metric Comparison}

\begin{lemma}[Metric Comparison]
\label{lem:metric_comparison}
If $\rho = \|z - z_0\|_z < 1$:
\begin{equation}
    (1-\rho)^2\, E \;\le\; E_0 \;\le\; (1+\rho)^2\, E.
\end{equation}
\end{lemma}

\section{Computable Step-Size Policy}

All quantities in the symmetric cone bounds are computable:
\begin{itemize}
    \item $E = \texttt{hessianNormSquared}(z, \text{target})$
    \item $E_0 = \texttt{hessianNormSquared}(z_0, \text{target})$
    \item $d_\infty^2$: recovered from
      $\alpha = \texttt{stepSize}(z, \text{target})$ via
      $d_\infty^2 = 2/\alpha$ when $\alpha < 1$
    \item $\sigma = E_0/E$
\end{itemize}

\begin{corollary}[Unified Frozen-J Step Size for Symmetric Cones]
\label{cor:sym_unified}
Define $\alpha = \texttt{stepSize}(z, \text{target})$.  The step size
\begin{equation}
    h_f = \min\!\left(\alpha,\; \frac{E_0}{2E + d_\infty^2 V}\right)
\end{equation}
guarantees descent provided $|\epsilon| \le E_0/2$.

Since $V$ is not directly available, we can use the conservative
substitution $V \le \nu/k^2$ (the gap bound) or simply use
\begin{equation}
    h_f = \min\!\left(\alpha,\; \frac{\sigma}{2 + d_\infty^2 \cdot \nu/k^2 / E}\right)
\end{equation}
where $\nu/k^2$ is the current barrier gap (computable).

In the near-field ($d_\infty^2 < 2$, so $\alpha = 1$): take $h_f = 1$
and check descent via $\frac{1}{2}d_\infty^2 + |1-\sigma| + |\epsilon|/V < 1$.
\end{corollary}

\subsection{Refactoring Condition}

\begin{proposition}[Refactoring Criterion]
\label{prop:refactor}
Take the frozen step if $\sigma = E_0/E \in [\tfrac{1}{4}, 4]$;
otherwise refactor.  This ensures $\rho < 1$ (within Dikin ellipsoids).
\end{proposition}


\end{document}
